\documentclass[11pt,a4paper]{article}
\usepackage[utf8]{inputenc}
\usepackage[T1]{fontenc}
\usepackage[french]{babel}
\usepackage{geometry}
\usepackage{booktabs}
\usepackage{longtable}
\usepackage{array}
\usepackage{xcolor}
\usepackage{colortbl}
\usepackage{graphicx}
\usepackage{fancyhdr}
\usepackage{titlesec}
\usepackage{tikz}
\usepackage{enumitem}
\usepackage{hyperref}
\usepackage{multicol}
\usepackage{listings}
\usepackage{tcolorbox}
\usepackage{amssymb}
\usepackage{float}
\usepackage{subcaption}
\usepackage{amsmath}
\usepackage{multirow}

\geometry{margin=2cm}
\pagestyle{fancy}
\fancyhf{}
\rhead{Atelier Machine Learning}
\lhead{Analyse Comportementale Clientèle}
\rfoot{Page \thepage}

% Couleurs personnalisées
\definecolor{headerblue}{RGB}{41, 65, 114}
\definecolor{lightgray}{RGB}{245, 245, 245}
\definecolor{accentorange}{RGB}{230, 126, 34}
\definecolor{codegreen}{RGB}{0,128,0}
\definecolor{codegray}{RGB}{128,128,128}
\definecolor{codepurple}{RGB}{128,0,128}
\definecolor{backcolour}{RGB}{248,248,248}
\definecolor{warningred}{RGB}{231, 76, 60}
\definecolor{successgreen}{RGB}{46, 204, 113}
\definecolor{infoBlue}{RGB}{52, 152, 219}
\definecolor{darkBlue}{RGB}{41, 65, 114}

\lstdefinestyle{mystyle}{
    backgroundcolor=\color{backcolour},   
    commentstyle=\color{codegreen},
    keywordstyle=\color{codepurple},
    numberstyle=\tiny\color{codegray},
    stringstyle=\color{codegreen},
    basicstyle=\ttfamily\footnotesize,
    breakatwhitespace=false,         
    breaklines=true,                 
    captionpos=b,                    
    keepspaces=true,                 
    numbers=left,                    
    numbersep=5pt,                  
    showspaces=false,                
    showstringspaces=false,
    showtabs=false,                  
    tabsize=2
}
\lstset{style=mystyle}

\tcbset{
    colback=backcolour,
    colframe=headerblue,
    fonttitle=\bfseries,
    boxrule=1pt,
    arc=3mm
}

\titleformat{\section}
{\color{headerblue}\normalfont\Large\bfseries}
{\thesection}{1em}{}[{\titlerule[2pt]}]

\titleformat{\subsection}
{\color{accentorange}\normalfont\large\bfseries}
{\thesubsection}{1em}{}

\titleformat{\subsubsection}
{\color{darkBlue}\normalfont\normalsize\bfseries}
{\thesubsubsection}{1em}{}

\newcolumntype{C}[1]{>{\centering\arraybackslash}p{#1}}
\newcolumntype{L}[1]{>{\raggedright\arraybackslash}p{#1}}

\begin{document}

\begin{titlepage}
\begin{tikzpicture}[remember picture,overlay]
    \fill[headerblue] (current page.north west) rectangle ([yshift=-5cm]current page.north east);
    \node[anchor=north west, xshift=2cm, yshift=-2cm, text=white] at (current page.north west) 
        {\Huge\textbf{Atelier Machine Learning}};
    \node[anchor=north west, xshift=2cm, yshift=-3.5cm, text=white] at (current page.north west) 
        {\Large\textit{Analyse Comportementale Clientèle Retail}};
\end{tikzpicture}

\vspace{6cm}

\begin{center}
{\LARGE\textbf{COMPTE RENDU}}\par
\vspace{1cm}
{\large Atelier Pratique : E-commerce de Cadeaux}\par
\vspace{2cm}
{\large \textbf{Préparé par Mariem Ben Abdallah}}
\vspace{3cm}

\begin{tikzpicture}
    \draw[headerblue, line width=2pt] (0,0) -- (10,0);
\end{tikzpicture}

\vspace{2cm}

{\large \textbf{Encadrante : Fadoua Drira}}

\vspace{5cm}

{\textbf{Année Universitaire : 2025-2026}}

\end{center}
\vfill
\end{titlepage}

\tableofcontents
\newpage

\section{Introduction et Contexte}

\subsection{Présentation du projet}
Ce projet s'inscrit dans le cadre de l'atelier pratique du module Machine Learning. 
L'objectif est de réaliser une \textbf{analyse comportementale complète} de la clientèle 
d'une entreprise e-commerce spécialisée dans la vente de cadeaux.

Le projet suit une méthodologie rigoureuse couvrant l'ensemble de la chaîne de traitement 
en data science : de l'exploration des données brutes jusqu'au déploiement d'une application 
web interactive.

\subsection{Objectifs métier}
L'entreprise e-commerce souhaite mieux comprendre sa clientèle pour :
\begin{itemize}[label={\color{headerblue}\textbullet}]
    \item \textbf{Personnaliser} ses stratégies marketing (offres ciblées, communications adaptées)
    \item \textbf{Réduire} le taux de départ des clients (churn) en anticipant les départs
    \item \textbf{Optimiser} son chiffre d'affaires en fidélisant les clients à forte valeur
    \item \textbf{Anticiper} les comportements à risque pour agir avant le départ
\end{itemize}

\subsection{Objectifs pédagogiques}
\begin{tcolorbox}[title=�� Compétences développées]
Ce projet couvre l'ensemble de la chaîne de traitement en data science :
\begin{multicols}{2}
\begin{itemize}
    \item Exploration des données
    \item Préparation et nettoyage
    \item Feature engineering
    \item Transformation (ACP)
    \item Modélisation (Clustering, Classification)
    \item Évaluation et interprétation
    \item Optimisation des hyperparamètres
    \item Déploiement (Flask)
\end{itemize}
\end{multicols}
\end{tcolorbox}

\subsection{Structure du rapport}
Ce rapport est organisé en plusieurs sections :
\begin{enumerate}
    \item \textbf{Exploration des données} : analyse statistique, valeurs manquantes, visualisations
    \item \textbf{Préparation des données} : nettoyage, feature engineering, encodage
    \item \textbf{Modélisation} : entraînement et comparaison de 4 modèles
    \item \textbf{Segmentation clientèle} : ACP et K-means pour identifier des profils
    \item \textbf{Déploiement} : application Flask et API REST
    \item \textbf{Conclusion} : bilan des résultats et perspectives
\end{enumerate}

\newpage
\section{Exploration et Analyse des Données}

\subsection{Présentation du dataset}
Le jeu de données fourni contient \textbf{4 372 clients} et \textbf{52 features}. 
Il s'agit de données intentionnellement imparfaites (valeurs manquantes, outliers, 
features redondantes) pour nous permettre de maîtriser l'ensemble des étapes de 
prétraitement.

\begin{tcolorbox}[title=�� Structure des données]
\begin{tabular}{L{4cm} C{3cm} C{3cm}}
\toprule
\textbf{Type de features} & \textbf{Nombre} & \textbf{Exemples} \\
\midrule
Numériques & 34 & Recency, Frequency, Monetary, Age \\
Catégorielles & 18 & RFMSegment, CustomerType, Region \\
\bottomrule
\end{tabular}
\end{tcolorbox}

\subsection{Analyse des types de données}
Nous avons commencé par analyser la structure du dataset :
\begin{itemize}
    \item \textbf{Features numériques} : variables quantitatives (RFM, quantités, ratios)
    \item \textbf{Features catégorielles} : segments clients, préférences, statuts
    \item \textbf{Features temporelles} : dates d'inscription, dernières connexions
    \item \textbf{Variable cible} : \texttt{Churn} (0 = fidèle, 1 = partant)
\end{itemize}

\subsection{Valeurs manquantes}
L'analyse des valeurs manquantes a révélé :

\begin{table}[H]
\centering
\begin{tabular}{L{5cm} C{2.5cm} C{2.5cm}}
\toprule
\textbf{Colonne} & \textbf{Manquantes} & \textbf{Pourcentage} \\
\midrule
Age & 1 311 & 29.99\% \\
AvgDaysBetweenPurchases & 79 & 1.81\% \\
\bottomrule
\end{tabular}
\caption{Valeurs manquantes par colonne}
\end{table}

\textbf{Interprétation :}
\begin{itemize}
    \item Les 79 clients sans délai entre achats sont ceux qui n'ont effectué \textbf{qu'un seul achat} (Frequency=1)
    \item L'âge manque pour environ 30\% des clients → nécessite une imputation
\end{itemize}

\subsection{Valeurs aberrantes et codes spéciaux}
Certaines colonnes contiennent des valeurs spéciales qui ne sont pas des erreurs mais 
des codes métier :

\begin{itemize}
    \item \textbf{SupportTicketsCount} : valeurs -1 (code pour "non applicable") et 999 (code pour "trop élevé")
    \item \textbf{SatisfactionScore} : valeurs -1, 0, 99 (codes spéciaux pour "non renseigné")
    \item \textbf{MonetaryTotal négatif} : 44 clients avec dépense totale négative (retours/annulations)
\end{itemize}

\subsection{Feature Engineering réalisé}
Pour enrichir notre jeu de données, nous avons créé de nouvelles features :

\begin{tcolorbox}[title=✨ Nouvelles features créées]
\begin{itemize}
    \item \texttt{HasNegativeMonetary} : flag binaire indiquant si le client a une dépense négative
    \item \texttt{RegYear, RegMonth, RegDay, RegWeekday} : extraction depuis la date d'inscription
    \item \texttt{IP\_IsPrivate} : indique si l'adresse IP est privée (1) ou publique (0)
    \item \texttt{IP\_Class} : classe de l'adresse IP (A, B, C, Other)
    \item \texttt{IP\_FirstOctet} : premier octet de l'adresse IP
\end{itemize}
\end{tcolorbox}

\subsection{Visualisations exploratoires}
Plusieurs visualisations ont été réalisées pour comprendre la distribution des données :

\begin{itemize}
    \item Distribution du MonetaryTotal (dépense totale)
    \item Corrélation entre les features numériques (heatmap)
    \item Taux de churn par mois d'inscription
    \item Répartition des clients par catégorie d'âge et de dépense
\end{itemize}

\newpage
\section{Préparation des Données}

\subsection{Nettoyage des données}
Le nettoyage a consisté à :
\begin{itemize}
    \item \textbf{Supprimer les colonnes constantes} : NewsletterSubscribed = "Yes" pour 100\% des clients
    \item \textbf{Parser les dates} : conversion de RegistrationDate en format datetime
    \item \textbf{Corriger les types} : s'assurer que chaque colonne a le bon type
    \item \textbf{Supprimer les identifiants} : CustomerID retiré (non prédictif)
\end{itemize}

\subsection{Suppression des features leakantes — Point critique}
Une attention particulière a été portée à l'identification et la suppression des 
features causant du \textbf{data leakage}. Le data leakage se produit lorsqu'une 
feature utilise une information qui ne serait pas disponible au moment de la prédiction.

\begin{tcolorbox}[title=⚠️ Features leakantes supprimées, colback=warningred!10]
\begin{itemize}
    \item \texttt{ChurnRiskCategory, RFMSegment, CustomerType, LoyaltyLevel} 
    → construites APRÈS avoir défini le churn (leakage direct)
    \item \texttt{FirstPurchaseDaysAgo, CustomerTenureDays} 
    → calculées sur la même période que le label (leakage temporel)
    \item \texttt{Recency} 
    → corrélation de 0.859 avec le churn (donnerait une précision artificielle de 100\%)
    \item \texttt{PreferredMonth, AvgDaysBetweenPurchases}
    → dérivées des mêmes informations temporelles
\end{itemize}
\end{tcolorbox}

\subsection{Gestion de la multicolinéarité}
La multicolinéarité désigne une forte corrélation entre plusieurs variables explicatives. 
Cela peut poser problème pour certains modèles (régression linéaire, logistique).

\textbf{Méthode utilisée :}
\begin{enumerate}
    \item Calcul de la matrice de corrélation
    \item Identification des paires avec |corrélation| > 0.8
    \item Suppression d'une feature par paire redondante
\end{enumerate}

\begin{table}[H]
\centering
\begin{tabular}{L{5cm} C{4cm} C{2cm}}
\toprule
\textbf{Feature 1 (conservée)} & \textbf{Feature 2 (supprimée)} & \textbf{Corrélation} \\
\midrule
Frequency & UniqueInvoices & 1.000 \\
UniqueProducts & UniqueDescriptions & 0.999 \\
MonetaryMin & MonetaryMax & -0.994 \\
MonetaryTotal & TotalQuantity & 0.922 \\
\bottomrule
\end{tabular}
\caption{Paires de features fortement corrélées}
\end{table}

\textbf{Résultat :} 7 features redondantes supprimées.

\subsection{Split Train/Test}
Le split des données est une étape cruciale pour évaluer correctement les performances 
d'un modèle sur des données jamais vues.

\begin{tcolorbox}[title=✅ Split réalisé]
\begin{itemize}
    \item Proportion : 80\% train / 20\% test
    \item Méthode : train\_test\_split avec random\_state=42 (reproductibilité)
    \item Stratification : sur la variable Churn (préserve la proportion des classes)
    \item Résultat : 3 497 clients en train, 875 en test
\end{itemize}
\end{tcolorbox}

\subsection{Imputation des valeurs manquantes — Correction critique}
L'imputation des valeurs manquantes doit être effectuée \textbf{APRÈS} le split pour 
éviter tout data leakage.

\begin{tcolorbox}[title=�� Bonne pratique respectée]
\begin{enumerate}
    \item Split Train/Test
    \item Identification des colonnes avec NaN dans X\_train
    \item SimpleImputer fit sur X\_train uniquement (stratégie = médiane)
    \item Transformation de X\_train et X\_test avec le même imputer
\end{enumerate}
\end{tcolorbox}

\textbf{Colonne imputée :} Age (1311 valeurs manquantes) → remplacées par la médiane.

\subsection{Encodage des variables catégorielles}
Les modèles de machine learning ne peuvent pas traiter directement les variables 
textuelles. Nous avons donc appliqué un \textbf{One-Hot Encoding}.

\begin{itemize}
    \item \textbf{12 variables catégorielles} encodées
    \item \textbf{drop\_first=True} pour éviter la multicolinéarité
    \item Nombre total de features après encodage : \textbf{103}
\end{itemize}

\textbf{Attention particulière :} Certaines catégories présentes uniquement dans 
le train (ex: Country\_Brazil, Region\_Amérique du Nord) ont été remplies avec 0 
dans le test.

\subsection{Normalisation des données}
Pour les modèles sensibles aux échelles (KNN, SVM), nous avons normalisé les données 
avec \textbf{StandardScaler}.

\begin{tcolorbox}[title=�� Normalisation après split]
\begin{itemize}
    \item StandardScaler fit sur X\_train\_encoded uniquement
    \item Transformation de X\_train\_encoded et X\_test\_encoded
    \item Sauvegarde du scaler pour l'application Flask
\end{itemize}
\end{tcolorbox}

\newpage
\section{Modélisation}

\subsection{Stratégie de modélisation}
Nous avons choisi de tester \textbf{4 modèles} aux architectures différentes pour 
comparer leurs performances et sélectionner le meilleur :

\begin{enumerate}
    \item \textbf{KNN} : modèle non paramétrique basé sur la proximité
    \item \textbf{Arbre de Décision} : modèle interprétable
    \item \textbf{Random Forest} : ensemble d'arbres (bagging)
    \item \textbf{Gradient Boosting} : ensemble séquentiel (boosting)
\end{enumerate}

\subsection{Modèle 1 : KNN (K-Nearest Neighbors)}

\subsubsection{Principe}
Le KNN classe un nouveau client en fonction des k clients les plus proches dans 
l'espace des features. C'est un modèle non paramétrique qui ne fait pas d'hypothèse 
sur la distribution des données.

\subsubsection{Optimisation réalisée}
Le KNN a bénéficié d'une optimisation poussée avec 3 stratégies :

\begin{enumerate}
    \item \textbf{SelectKBest} : sélection des 20 meilleures features (test de Fisher)
    \item \textbf{PCA} : réduction de dimension (testée mais moins performante)
    \item \textbf{SelectKBest + PCA} : combinaison des deux approches
\end{enumerate}

\textbf{Paramètres optimaux trouvés :}
\begin{itemize}
    \item Nombre de voisins (k) : 15
    \item Poids : distance (les voisins proches ont plus d'influence)
    \item Métrique : Manhattan
    \item Features sélectionnées : 20 meilleures
\end{itemize}

\subsubsection{Résultats KNN}
\begin{table}[H]
\centering
\begin{tabular}{L{3cm} C{3cm} C{3cm} C{2cm}}
\toprule
\textbf{Métrique} & \textbf{KNN Base} & \textbf{KNN Optimisé} & \textbf{Gain} \\
\midrule
Accuracy & 0.753 & 0.878 & +12.5\% \\
Recall & 0.430 & 0.825 & +39.5\% \\
F1-Score & 0.536 & 0.818 & +28.1\% \\
ROC-AUC & 0.789 & 0.925 & +13.6\% \\
\bottomrule
\end{tabular}
\caption{Amélioration du KNN après optimisation}
\end{table}

\subsection{Modèle 2 : Arbre de Décision}

\subsubsection{Principe}
L'arbre de décision segmente l'espace des features en régions homogènes selon la 
variable cible. Il est très interprétable mais sujet au sur-apprentissage.

\subsubsection{Optimisation réalisée}
\begin{itemize}
    \item Critère : entropie
    \item Contrainte : max\_depth=5 (pour éviter le sur-apprentissage)
    \item Class\_weight = 'balanced' (gestion du déséquilibre)
    \item Min\_samples\_split = 20, min\_samples\_leaf = 10
\end{itemize}

\begin{tcolorbox}[title=�� Impact de la profondeur maximale]
\begin{tabular}{L{4cm} C{3cm} C{3cm}}
\toprule
& \textbf{Sans contrainte} & \textbf{max\_depth=5} \\
\midrule
Profondeur & 18 & 5 \\
Accuracy Train & 1.000 & 0.874 \\
Accuracy Test & 0.839 & 0.864 \\
Écart & 16.1\% & 1.0\% \\
\bottomrule
\end{tabular}
\end{tcolorbox}

\subsubsection{Résultats Arbre de Décision}
\begin{itemize}
    \item \textbf{F1-Score} : 0.819
    \item \textbf{Recall} : 0.924 (le meilleur parmi tous les modèles !)
    \item \textbf{ROC-AUC} : 0.944
\end{itemize}

L'arbre de décision détecte \textbf{92.4\% des clients partants} — excellent pour 
une stratégie de fidélisation préventive.

\subsection{Modèle 3 : Random Forest}

\subsubsection{Principe}
La Random Forest construit de nombreux arbres de décision sur des échantillons 
bootstrappés et agrège leurs prédictions. Cela réduit la variance et améliore 
la généralisation.

\subsubsection{Paramètres utilisés}
\begin{itemize}
    \item n\_estimators = 200 arbres
    \item max\_depth = 10
    \item min\_samples\_split = 10
    \item max\_features = 'sqrt'
    \item class\_weight = 'balanced'
    \item oob\_score = True (validation out-of-bag)
\end{itemize}

\subsubsection{Résultats Random Forest}
\begin{itemize}
    \item \textbf{F1-Score} : 0.783
    \item \textbf{ROC-AUC} : 0.907
    \item \textbf{OOB Score} : 0.845
\end{itemize}

\subsection{Modèle 4 : Gradient Boosting — Meilleur modèle}

\subsubsection{Principe}
Le Gradient Boosting construit les arbres séquentiellement : chaque nouvel arbre 
corrige les erreurs du précédent. C'est souvent le modèle le plus performant sur 
les données tabulaires.

\subsubsection{Optimisation par GridSearchCV}
Nous avons effectué une recherche d'hyperparamètres sur 3 dimensions :

\begin{table}[H]
\centering
\begin{tabular}{L{4cm} C{3cm} C{3cm}}
\toprule
\textbf{Paramètre} & \textbf{Valeurs testées} & \textbf{Meilleure valeur} \\
\midrule
n\_estimators & 100, 200 & 200 \\
learning\_rate & 0.01, 0.1, 0.2 & 0.01 \\
max\_depth & 3, 5, 7 & 5 \\
\bottomrule
\end{tabular}
\caption{GridSearchCV - Gradient Boosting}
\end{table}

\textbf{Validation croisée :} F1 moyen = 0.827 (5 folds)

\subsubsection{Résultats Gradient Boosting}
\begin{tcolorbox}[title=�� Modèle final retenu, colback=successgreen!10]
\begin{itemize}
    \item \textbf{F1-Score} : \textbf{0.8297}
    \item \textbf{ROC-AUC} : \textbf{0.9471}
    \item \textbf{Accuracy} : 89.3\%
    \item \textbf{Precision} : 0.877
    \item \textbf{Recall} : 0.787
\end{itemize}
\end{tcolorbox}

\subsection{Top features importantes}
L'analyse des importances des features du Gradient Boosting révèle :

\begin{table}[H]
\centering
\begin{tabular}{L{5cm} C{3cm} C{3cm}}
\toprule
\textbf{Feature} & \textbf{Importance} & \textbf{Interprétation} \\
\midrule
FavoriteSeason\_Hiver & 22.4\% & Saison d'achat préférée \\
FavoriteSeason\_Été & 20.0\% & Saison d'achat préférée \\
Frequency & 19.6\% & Fréquence d'achat \\
FavoriteSeason\_Printemps & 14.2\% & Saison d'achat préférée \\
TotalTransactions & 9.0\% & Nombre total de transactions \\
MonetaryTotal & 3.3\% & Dépense totale \\
\bottomrule
\end{tabular}
\caption{Top 6 features importantes}
\end{table}

\begin{tcolorbox}[title=�� Analyse des features importantes, colback=infoBlue!10]
\textbf{Observation :} Les features FavoriteSeason (Hiver, Été, Printemps) représentent 
56.5\% de l'importance totale du modèle.

\textbf{Explication possible :} Les données couvrent une période spécifique (2009-2011) 
où la saison d'achat est corrélée avec la période de labellisation du churn. 
Ce n'est pas nécessairement du leakage mais une corrélation spurieuse à mentionner 
dans les limites du projet.
\end{tcolorbox}

\subsection{Tableau comparatif final}

\begin{table}[H]
\centering
\begin{tabular}{L{4cm} C{2.5cm} C{2.5cm} C{2.5cm} C{2.5cm}}
\toprule
\textbf{Modèle} & \textbf{Accuracy} & \textbf{Recall} & \textbf{F1-Score} & \textbf{ROC-AUC} \\
\midrule
\textbf{Gradient Boosting} & \textbf{0.8926} & 0.7869 & \textbf{0.8297} & \textbf{0.9471} \\
Arbre Décision & 0.8640 & \textbf{0.9244} & 0.8189 & 0.9440 \\
KNN Optimisé & 0.8777 & 0.8247 & 0.8177 & 0.9249 \\
Random Forest & 0.8549 & 0.7869 & 0.7829 & 0.9071 \\
\bottomrule
\end{tabular}
\caption{Comparaison des performances des modèles}
\end{table}

\subsection{Choix du modèle final}
Nous retenons le \textbf{Gradient Boosting} comme modèle final pour les raisons suivantes :

\begin{enumerate}
    \item \textbf{Meilleur F1-Score} (0.83) : équilibre optimal entre précision et rappel
    \item \textbf{Meilleur ROC-AUC} (0.947) : excellent pouvoir discriminant
    \item \textbf{Meilleure Accuracy} (89.3\%) : plus de prédictions correctes
    \item \textbf{Modèle robuste} : bien adapté aux données tabulaires complexes
\end{enumerate}

\newpage
\section{Segmentation Clientèle (ACP + K-means)}

\subsection{Objectif de la segmentation}
La segmentation clientèle vise à identifier des groupes homogènes de clients pour 
adapter les stratégies marketing. Nous utilisons une approche non supervisée 
(ACP + K-means).

\subsection{Prétraitement pour la segmentation}
\begin{enumerate}
    \item \textbf{Sélection des features numériques} : 27 colonnes conservées
    \item \textbf{Imputation} : médiane pour les valeurs manquantes
    \item \textbf{Normalisation} : StandardScaler (indispensable pour K-means)
    \item \textbf{Retrait des outliers} : Z-score > 3 sur au moins une dimension
\end{enumerate}

\begin{tcolorbox}[title=�� Outliers détectés]
\begin{itemize}
    \item \textbf{Outliers identifiés} : 817 clients (23.4\% du dataset)
    \item \textbf{Clients conservés} : 2 680 clients (76.6\%)
    \item \textbf{Profil des outliers} : fréquence moyenne 7.7 achats, dépense moyenne 3842€
\end{itemize}
\end{tcolorbox}

\subsection{Analyse en Composantes Principales (ACP)}

\subsubsection{Principe}
L'ACP est une technique de réduction de dimension qui transforme les variables 
corrélées en variables indépendantes (composantes principales) tout en conservant 
le maximum de variance.

\subsubsection{Résultats de l'ACP}
\begin{itemize}
    \item \textbf{14 composantes} retenues pour expliquer 80\% de la variance
    \item \textbf{20 composantes} nécessaires pour 95\% de la variance
    \item Visualisation 2D possible pour l'interprétation
\end{itemize}

\subsection{Détermination du nombre de clusters (K-means)}

\subsubsection{Principe de K-means}
K-means partitionne les données en k clusters en minimisant la distance intra-cluster. 
Chaque cluster est représenté par son centroïde.

\subsubsection{Méthode du coude}
Nous avons utilisé la méthode du coude (Elbow method) qui consiste à tracer l'inertie 
en fonction du nombre de clusters k. Le "coude" indique le k optimal.

\begin{itemize}
    \item Test de k=2 à 10
    \item Coude détecté autour de k=3 ou k=4
    \item Analyse des tailles minimales : k=4 retenu (plus petit cluster = 222 clients, 8.3\%)
\end{itemize}

\subsection{Résultats de la segmentation}

\begin{table}[H]
\centering
\begin{tabular}{C{2cm} L{4cm} C{2cm} C{2cm} C{2cm}}
\toprule
\textbf{Cluster} & \textbf{Nom du segment} & \textbf{Taille} & \textbf{Churn} & \textbf{Profil} \\
\midrule
0 & Clients Réguliers & 222 (8.3\%) & 19.4\% & 6.1 achats, 1735£ \\
1 & Occasionnels à Risque & 744 (27.8\%) & 29.7\% & 3.5 achats, 1121£ \\
2 & Occasionnels à Risque & 747 (27.9\%) & 34.3\% & 2.9 achats, 902£ \\
3 & Clients Réguliers & 967 (36.1\%) & 37.6\% & 5.7 achats, 1697£ \\
- & Outliers & 817 (23.4\%) & 34.1\% & Analyse individuelle \\
\bottomrule
\end{tabular}
\caption{Récapitulatif des segments clients}
\end{table}

\subsection{Interprétation des clusters}

\subsubsection{Cluster 0 — Clients Réguliers (Risque faible)}
\begin{itemize}
    \item \textbf{Taille} : 222 clients (8.3\%)
    \item \textbf{Taux de churn} : 19.4\% (le plus bas)
    \item \textbf{Profil} : fréquence d'achat élevée (6.1), dépense importante (1735£)
    \item \textbf{Action recommandée} : programme de fidélité, upselling, cross-selling
\end{itemize}

\subsubsection{Cluster 1 — Occasionnels à Risque modéré}
\begin{itemize}
    \item \textbf{Taille} : 744 clients (27.8\%)
    \item \textbf{Taux de churn} : 29.7\%
    \item \textbf{Profil} : fréquence moyenne (3.5), dépense modérée (1121£)
    \item \textbf{Action recommandée} : newsletter ciblée, promotions personnalisées
\end{itemize}

\subsubsection{Cluster 2 — Occasionnels à Risque élevé}
\begin{itemize}
    \item \textbf{Taille} : 747 clients (27.9\%)
    \item \textbf{Taux de churn} : 34.3\%
    \item \textbf{Profil} : fréquence faible (2.9), dépense faible (902£)
    \item \textbf{Action recommandée} : offres de réactivation urgentes
\end{itemize}

\subsubsection{Cluster 3 — Clients Réguliers (Risque élevé)}
\begin{itemize}
    \item \textbf{Taille} : 967 clients (36.1\%)
    \item \textbf{Taux de churn} : 37.6\% (le plus élevé des clusters)
    \item \textbf{Profil} : fréquence élevée (5.7) mais churn élevé (paradoxe)
    \item \textbf{Action recommandée} : campagne anti-churn prioritaire
\end{itemize}

\subsubsection{Outliers — Cas atypiques}
\begin{itemize}
    \item \textbf{Taille} : 817 clients (23.4\%)
    \item \textbf{Profil} : comportements extrêmes sur au moins une dimension
    \item \textbf{Action recommandée} : analyse individuelle (pas de campagne marketing générique)
\end{itemize}

\newpage
\section{Déploiement avec Flask}

\subsection{Choix technologique}
Flask a été choisi pour le déploiement car c'est un micro-framework Python :
\begin{itemize}
    \item Léger et flexible
    \item Facile à intégrer avec Scikit-learn
    \item Idéal pour créer des API REST
    \item Parfait pour un projet académique
\end{itemize}

\subsection{Architecture de l'application}
L'application Flask se compose de :
\begin{enumerate}
    \item \textbf{app.py} : le serveur backend
    \item \textbf{templates/index.html} : l'interface utilisateur
    \item \textbf{Modèles sauvegardés} : best\_model\_churn.pkl, scaler.pkl
\end{enumerate}

\subsection{API REST implémentée}
\begin{tcolorbox}[title=�� Endpoints disponibles]
\begin{tabular}{L{5cm} L{5cm}}
\toprule
\textbf{Endpoint} & \textbf{Fonction} \\
\midrule
GET / & Interface web \\
POST /predict & Prédiction via formulaire \\
POST /api/predict & Prédiction via API JSON \\
GET /health & Vérification de santé \\
GET /debug & Informations de debug \\
\bottomrule
\end{tabular}
\end{tcolorbox}

\subsection{Fonctionnalités implémentées}

\subsubsection{Interface utilisateur}
\begin{itemize}
    \item Formulaire avec 6 champs (Recency, Frequency, Monetary, Age, Satisfaction, Tickets)
    \item Boutons d'exemples pré-remplis (client fidèle / client à risque)
    \item Jauge de risque visuelle (couleurs vert/orange/rouge)
    \item Scores RFM calculés automatiquement
    \item Conseils personnalisés selon le niveau de risque
    \item Historique des 10 dernières prédictions
\end{itemize}

\subsubsection{Prétraitement en temps réel}
L'application reconstruit un vecteur de features complet (103 dimensions) à partir 
des données saisies :
\begin{itemize}
    \item Valeurs numériques saisies par l'utilisateur
    \item Valeurs par défaut pour les features non renseignées
    \item Encodage one-hot des variables catégorielles
    \item Normalisation avec le scaler sauvegardé
\end{itemize}

\subsection{Exemples de prédictions}

\begin{table}[H]
\centering
\begin{tabular}{C{3cm} C{2cm} C{2cm} C{2cm} C{2.5cm} C{2cm}}
\toprule
\textbf{Profil} & \textbf{Freq} & \textbf{Monetary} & \textbf{Sat} & \textbf{Risque} & \textbf{Pred} \\
\midrule
Client Premium & 50 & 10000€ & 5 & 15\% & Fidèle \\
Client Moyen & 5 & 500€ & 4 & 40\% & Fidèle \\
Client Critique & 1 & 50€ & 1 & 89\% & Partant \\
\bottomrule
\end{tabular}
\caption{Exemples de prédictions de l'interface}
\end{table}

\subsection{Logs et monitoring}
L'application affiche des logs détaillés dans la console :
\begin{itemize}
    \item Données reçues du formulaire
    \item Encodage des variables catégorielles
    \item Vérification du vecteur final (NaN, shape)
    \item Résultat de la prédiction (probabilité, niveau)
\end{itemize}

\subsection{Utilisation de l'API REST}
\begin{lstlisting}[language=bash]
# Exemple d'appel API avec curl
curl -X POST http://127.0.0.1:5000/api/predict \
  -H "Content-Type: application/json" \
  -d '{
    "recency": 30,
    "frequency": 5,
    "monetary": 500,
    "age": 35,
    "satisfaction": 4,
    "support_tickets": 0,
    "favorite_season": "Hiver",
    "spending_category": "Medium",
    "account_status": "Active"
  }'
\end{lstlisting}

\newpage
\section{Conclusion et Perspectives}

\subsection{Résumé des réalisations}
\begin{tcolorbox}[title=✅ Bilan du projet]
\begin{itemize}
    \item \textbf{Exploration} : analyse complète de 52 features, identification des anomalies
    \item \textbf{Préparation} : nettoyage, feature engineering, suppression leakage (9 features)
    \item \textbf{Modélisation} : 4 modèles testés, Gradient Boosting retenu (F1=0.83, AUC=0.947)
    \item \textbf{Segmentation} : 4 clusters clients + outliers, recommandations actionnables
    \item \textbf{Déploiement} : application Flask fonctionnelle avec API REST
\end{itemize}
\end{tcolorbox}

\subsection{Points forts du projet}
\begin{enumerate}
    \item \textbf{Rigueur scientifique} : split avant imputation, pas de data leakage
    \item \textbf{Modèle performant} : F1-Score de 0.83, AUC de 0.947
    \item \textbf{Segmentation exploitable} : 4 clusters aux profils distincts
    \item \textbf{Interface professionnelle} : jauge de risque, scores RFM, historique
    \item \textbf{Code versionné} : GitHub avec commits réguliers
\end{enumerate}

\subsection{Limites identifiées}
\begin{itemize}
    \item \textbf{FavoriteSeason domine le modèle} (56.5\% de l'importance) → possible biais temporel
    \item \textbf{23.4\% d'outliers} retirés de la segmentation → perte d'information
    \item \textbf{Donées datées} : période 2009-2011, centrées sur le Royaume-Uni
    \item \textbf{Age imputé par médiane} : 30\% de valeurs manquantes
    \item \textbf{KNN moins performant} sur données tabulaires complexes (F1=0.818)
\end{itemize}

\subsection{Perspectives d'amélioration}
\begin{multicols}{2}
\begin{itemize}
    \item Test de modèles deep learning (MLP, TabNet)
    \item Ajout de données plus récentes (2020-2025)
    \item Intégration de SHAP pour l'interprétabilité
    \item Mise en production sur cloud (AWS, GCP, Heroku)
    \item Dashboard temps réel avec Plotly Dash
    \item A/B testing des recommandations marketing
    \item Système de monitoring du modèle (drift detection)
    \item Ajout de plus de features externes (météo, économie)
\end{itemize}
\end{multicols}

\subsection{Livrables du projet}
\begin{itemize}
    \item \textbf{Code source} : dépôt GitHub complet
    \item \textbf{Application Flask} : interface de prédiction fonctionnelle
    \item \textbf{Modèle entraîné} : best\_model\_churn.pkl
    \item \textbf{Notebooks} : analyses, modélisation, segmentation
    \item \textbf{Rapport} : présent document (LaTeX)
    \item \textbf{Présentation} : slides pour la soutenance
\end{itemize}

\newpage
\section*{Annexes}
\addcontentsline{toc}{section}{Annexes}

\subsection*{A. Structure complète du projet}
\begin{verbatim}
ProjetML-1/
├── app/
│   ├── app.py                 # Application Flask
│   └── templates/
│       └── index.html         # Interface utilisateur
├── models/
│   ├── best_model_churn.pkl   # Modèle Gradient Boosting
│   ├── scaler.pkl             # StandardScaler
│   ├── imputer.pkl            # SimpleImputer
│   ├── feature_names.pkl      # Noms des 103 features
│   └── ...                    # Autres modèles (kmeans, pca)
├── notebooks/
│   ├── test.py                # Script principal
│   └── ...                    # Notebooks d'analyse
├── data/
│   ├── raw/                   # Données brutes
│   ├── processed/             # Données nettoyées
│   └── train_test/            # Split train/test
├── reports/
│   ├── segmentation_dashboard.png
│   ├── arbre_decision.png
│   └── segmentation_recommandations.csv
├── src/
│   ├── preprocessing.py
│   ├── train_model.py
│   └── utils.py
├── requirements.txt
├── README.md
└── .gitignore
\end{verbatim}

\subsection*{B. Environnement technique}
\begin{table}[H]
\centering
\begin{tabular}{L{4cm} L{6cm}}
\toprule
\textbf{Catégorie} & \textbf{Technologies} \\
\midrule
Langage & Python 3.10+ \\
Data & Pandas, NumPy \\
Visualisation & Matplotlib, Seaborn \\
Machine Learning & Scikit-learn (GB, RF, KNN, DecisionTree) \\
Déploiement & Flask, HTML/CSS \\
Versioning & Git, GitHub \\
\bottomrule
\end{tabular}
\end{table}

\subsection*{C. Installation et exécution}
\begin{lstlisting}[language=bash]
# Cloner le repository
git clone https://github.com/myriamBenAbd0607/ProjetML.git
cd ProjetML

# Créer l'environnement virtuel
python -m venv venv
source venv/bin/activate  # Linux/Mac
venv\Scripts\activate     # Windows

# Installer les dépendances
pip install -r requirements.txt

# Lancer l'application Flask
cd app
python app.py

# Ouvrir dans le navigateur
# http://127.0.0.1:5000
\end{lstlisting}

\subsection*{D. Références}
\begin{itemize}
    \item Documentation Scikit-learn : \url{https://scikit-learn.org}
    \item Documentation Flask : \url{https://flask.palletsprojects.com}
    \item Documentation Pandas : \url{https://pandas.pydata.org}
    \item Support de cours : Module Machine Learning - GI2
\end{itemize}

\end{document}